\chapter{Conclusion}
\label{ch:conclusion}

\section{Objectives revisited}

Table~\ref{tab:objectives} maps each objective from Section~\ref{sec:objectives}
to the evidence bearing on it. Five were met outright; one was met, but in a form
different from the one originally proposed.

\begin{table}[htbp]
\centering
\footnotesize
\caption[Objectives against evidence]{Objectives against the evidence in this
report.}
\label{tab:objectives}
\begin{tabular}{@{}clp{5.5cm}l@{}}
\toprule
& \textbf{Objective} & \textbf{Evidence} & \textbf{Status} \\
\midrule
O1 & Three agents + summariser & \S\ref{sec:grounding}, Table~\ref{tab:agent-config};
     interests traced to Doc 1 categories & Met \\
O2 & Deterministic orchestration, validated output & \S\ref{sec:orchestration};
     structural compliance 0.9986 over 2{,}865 responses & Met \\
O3 & Conditionally verifying court agent & \S\ref{sec:court-design};
     0 of 9 flagged instances Court-originated & Met \\
O4 & Real documents to structured scenarios & \S\ref{sec:extraction};
     23 cases ingested through the user upload path & Met \\
O5 & Verified corpus, baselines, ablations & \S\ref{sec:baselines}--\S\ref{sec:ablation};
     13 of 14 merits dispositions matched & Met \\
O6 & Something the client can use & \S\ref{sec:risk-screen};
     rule-based screen, not the simulator & Met, reframed \\
\bottomrule
\end{tabular}
\end{table}

\section{What the research questions turned out to have as answers}

\textbf{RQ1, on prompt engineering without fine-tuning.} Yes broadly, with a
caveat that only surfaced under testing. Agents specified purely through role
prompts and injected context do sustain distinct, role-appropriate positions:
the authority concedes rarely on opening and about half the time thereafter, the
bidder concedes far less, and each reports the opposite kind of outcome against
its own declared fallback. The caveat is that a prompt developed against one
dispute shape doesn't automatically generalise to another. Confronted with a
case where nobody was ever actually scored, both negotiating agents invented a
scoring dispute anyway, simply because that was the only script either of them
had.

\textbf{RQ2, on which part of the protocol is load-bearing.} The adjudicator,
and essentially nothing else. Removing it deadlocks every single run. In
hindsight that result is close to trivial, since nothing else in the graph can
set the resolution flag, but it remains the one unambiguous architectural
finding in this
project, and it's worth more than the accuracy figures precisely because the
accuracy figures don't distinguish this system from a single prompt.

\textbf{RQ3, on constraining an LLM adjudicator's arithmetic.} Largely yes, and
the boundary of that discipline is now mapped rather than merely assumed. Gating
exact computation on the presence of both a formula and every input it needs,
and banning unstated numbers even when hedged as illustrative, held across 425
qualitative assessments with no Court-originated fabrication among the nine
flagged. It doesn't hold universally, though: there is a documented instance of
the Court inventing an entire illustrative dataset on a case it had
misclassified as numeric, and the arithmetic defect found on the client's own
scoring table resisted three separate prompt revisions. The discipline
constrains what the adjudicator computes; it does not constrain what premise it
is computing on.

\textbf{RQ4, on evaluating in the absence of ground truth.} What this project
arrived at is a stack of partial checks with explicit sample sizes attached, and
a deliberate refusal to collapse them into a single score. The stack needs an
architecture check, because a system can agree with reality for reasons that
have nothing to do with its design. It needs fabrication screens, because a
correct conclusion reached from invented evidence looks identical from the
outside to a correct conclusion reached honestly. And it needs to report
negative results, since the two most informative findings in this whole
project, the zero-shot tie and the readability shortfall, both count against the
artefact rather than for it.

\section{Reflection on the approach}

The scope ended up wide: five agent types, a state-machine orchestrator, a
document ingestion path, a streaming API, a client-facing web application, a
rule-based risk screen, a 23-case verified legal corpus, and eight analysis
scripts, all held together by 92 tests. The complexity that actually mattered,
though, was never any of those pieces individually. It was that the domain
punishes exactly the thing generative models are best at. A language model asked
about a procurement dispute will produce something fluent, structured and
legally shaped and the harder the case, the more confident it tends to sound.
Almost every engineering decision in Chapter~\ref{ch:methodology} traces back to
that one fact: schema-validated output so malformed responses can't pass
silently, per-response compliance counters so repairs stay visible instead of
disappearing into the log, an extraction agent instructed as firmly against
dropping real numbers as against inventing new ones, a minimum source-length
guard because an empty PDF once produced a complete fictional dispute out of
nothing, and a conditional verification protocol whose whole purpose is getting
the adjudicator to say ``I cannot compute this'' precisely when it can't.

There are two things I would do differently given the chance again. First, I
would build the corpus before building the agents. The synthetic scenarios
written early on were AI-authored, which quietly made every result derived from
them circular, and they were eventually deleted along with seven directories of
evaluation output. That cost roughly a month, though not for nothing: the reason
those results had to go is itself a methodological finding, and it is the reason
the corpus in this report is 23 verified real judgments rather than a set of
plausible fictions. A model-authored scenario cannot falsify a model-based
system, and discovering that early enough to rebuild on real case law is what
makes every figure in Chapter~\ref{ch:evaluation} mean anything.
Second, I would have started testing on documents I hadn't chosen myself much
sooner. Two of the most informative defects in this project came out of a
single hour spent running a Fusion21 employee's own files, and both sat in code
paths that months of my own test cases had never once exercised, because
without quite noticing it I had been selecting inputs shaped like the failures
I already understood.

On execution quality, the structural layer is strong and the measurements say so:
0.9986 clean structured responses across 2{,}865 calls, a green 92-case test
suite, every figure in this report reproducible from committed logs by a named
command, and provenance carried all the way through to the analytics interface so
a partial batch can never be mistaken for a finished one. The reasoning layer is
where the limits sit, and the important point is that they are limits this
project can state numerically rather than suspect: a 23.1\% wrong-regime citation
rate, a summariser that fails its own accessibility claim by 35 points of Flesch
Reading Ease, an outcome vocabulary that cannot express three of the real remedies
present in its own corpus, and one arithmetic defect that prompt engineering alone
never managed to close. Each of those is quantified, traced to a named case or
command, and paired with the specific fix it needs in \S\ref{sec:future}. A
system whose failure modes are this precisely bounded is a more useful starting
point than one whose failures have simply not been looked for.

\section{Future work}
\label{sec:future}

Four directions follow directly from findings already reported above, rather
than from speculation about what might be nice to try.

\textbf{Ground the citations.} A 23.1\% wrong-regime rate is the strongest
argument anywhere in this report for the retrieval layer that was scoped out
earlier. Retrieval over the Procurement Act 2023 and the TCC guidance targets
exactly the failure observed here: the model simply does not know which regime
governs a given year, and no instruction can hand it that knowledge the way
retrieval actually can.

\textbf{Check premises the way arithmetic is checked.} The Court agent verifies
numbers against the scenario record but accepts factual assertions from the
negotiating agents without ever checking them against that same record. The
Faraday result shows exactly what that costs. A grounding check applied to
negotiating-agent statements, built on the same machinery the numeric screen
already provides, is a concrete and well-evidenced next step rather than an
open-ended one.

\textbf{Extend the remedy vocabulary beyond scoring challenges.} The corpus
itself already contains three real remedies the system simply cannot express: a
declaration of ineffectiveness, financial penalties on the authority, and a
court-determined re-ranking. Each of these is a documented gap tied to a named
case, not a hypothetical concern.

\textbf{Calibrate the risk screen against real outcomes.} This is the only thing
this project actually asks of Fusion21, and it's the highest-value item on this
whole list. Anonymised pre-action data from their member base (the same profile
fields the screen already requests, plus whether a challenge ended up following)
would let a severity-weighted heuristic be replaced or validated by a properly
fitted model, and would let the risk score be reported as an actual probability
rather than just an ordering. It is also the only component in this system whose
central claim is directly falsifiable, which is exactly why it deserves the
data most.

The broader conclusion here is narrower than the one I expected to end up
writing. A multi-agent LLM system can conduct a procedurally coherent
procurement dispute negotiation, and it can be constrained not to invent the
arithmetic it reasons from. What it cannot yet be shown to do is reason better
than a single well-posed prompt, and it can still be led into confident nonsense
by a premise it was simply handed rather than one it worked out itself. The
value in this work sits less in the simulator than in the evidence about what
such a system does when nobody is checking, and in the demonstration that
finding that out requires screens built specifically to catch a system that is
wrong in a fluent and entirely plausible way.